%%%%%%%%%%%%%%%%%%%%%%%%%%%%%%%%%%%%%%%%%
% Developer CV
% LaTeX Class
% Version 2.0 (12/10/23)
%
% This class originates from:
% http://www.LaTeXTemplates.com
%
% Authors:
% Omar Roldan
% Based on a template by  Jan Vorisek (jan@vorisek.me)
% Based on a template by Jan Küster (info@jankuester.com)
% Modified for LaTeX Templates by Vel (vel@LaTeXTemplates.com)
%
% License:
% The MIT License (see included LICENSE file)
%
%%%%%%%%%%%%%%%%%%%%%%%%%%%%%%%%%%%%%%%%%

%----------------------------------------------------------------------------------------
%	PACKAGES AND OTHER DOCUMENT CONFIGURATIONS
%----------------------------------------------------------------------------------------

\documentclass[9pt]{developercv} % Default font size, values from 8-12pt are recommended
\usepackage{multicol}
\setlength{\columnsep}{0mm}
%----------------------------------------------------------------------------------------

\begin{document}

%----------------------------------------------------------------------------------------
%	TITLE AND CONTACT INFORMATION
%----------------------------------------------------------------------------------------

\begin{minipage}[t]{0.5\textwidth}
	\vspace{-\baselineskip} % Required for vertically aligning minipages

	{ \fontsize{16}{20} \textcolor{black}{\textbf{\MakeUppercase{Yannick Lang}}}} % First name

	\vspace{6pt}

	{\Large M.Sc. Angewandte Informatik} % Career or current job title
\end{minipage}
\hfill
\begin{minipage}[t]{0.2\textwidth} % 20% of the page width for the first row of icons
	\vspace{-\baselineskip} % Required for vertically aligning minipages

	% The first parameter is the FontAwesome icon name, the second is the box size and the third is the text
	\icon{Globe}{11}{\href{https://y-lang.eu}{y-lang.eu}}\\
	\icon{Mobile}{11}{0151 70008124}\\
	\icon{MapMarker}{11}{Bamberg, DE}\\

\end{minipage}
\begin{minipage}[t]{0.27\textwidth} % 27% of the page width for the second row of icons
	\vspace{-\baselineskip} % Required for vertically aligning minipages

	\icon{Envelope}{11}{\href{mailto:mail@y-lang.eu}{mail@y-lang.eu}}\\
	\icon{Github}{11}{\href{https://github.com/layaxx}{github.com/layaxx}}\\
	\icon{LinkedinSquare}{11}{\href{https://www.linkedin.com/in/y-lang/}{/in/y-lang}}\\

\end{minipage}


%----------------------------------------------------------------------------------------
%	INTRODUCTION, SKILLS AND TECHNOLOGIES
%----------------------------------------------------------------------------------------

\begin{minipage}[t]{0.46\textwidth}
	\cvsect{Übersicht}
	\vspace{-6pt}


	Master-Absolvent der Angewandten Informatik mit Fokus auf Full-Stack-Webentwicklung (Next.js, Python) und Machine Learning. Erfahrung in der Leitung von IT-Teams, Serveradministration und der Implementierung moderner Softwarelösungen im Unternehmens- und Vereinskontext.
\end{minipage}
\hfill % Whitespace between
\begin{minipage}[t]{0.465\textwidth}
	\cvsect{Skills}
	\vspace{-6pt}

	\begin{minipage}[t]{0.28\textwidth}
		\textbf{Languages:}
	\end{minipage}
	\hfill
	\begin{minipage}[t]{0.71\textwidth}
		Javascript/Typescript, Python, Java, SQL, R, C++, C.
	\end{minipage}
	\vspace{3mm}

	\begin{minipage}[t]{0.28\textwidth}
		\textbf{Technologies:}
	\end{minipage}
	\hfill
	\begin{minipage}[t]{0.71\textwidth}
		React, Next.js, Docker, CI/CD, PyTorch.
	\end{minipage}
	\vspace{-1mm}

	\begin{minipage}[t]{0.28\textwidth}
		\textbf{Principles:}
	\end{minipage}
	\hfill
	\begin{minipage}[t]{0.71\textwidth}
		Clean Code, Test-Driven Development, Code Reviews, Git best practices.
	\end{minipage}

\end{minipage}

%----------------------------------------------------------------------------------------
%	Projects
%----------------------------------------------------------------------------------------
\cvsect{Erfahrung}
\begin{entrylist}
	\entry
	{seit 9/2024}
	{Webentwickler (Teilzeit)}
	{LivingLogic AG}
	{\vspace{-10pt}
		\begin{itemize}[noitemsep,topsep=1pt,parsep=1pt,partopsep=0pt, leftmargin=-1pt]
			\item Webentwicklung im Kontext der Low-Code-Plattform LivingApps
			\item Entwicklung neuer Features in bestehenden Applikationen durch Frontend- und Backend-Anpassungen
			\item Integration von Large Language Models zur Unterstützung von Anwendern
		\end{itemize}
		\texttt{Javascript} \slashsep \texttt{Python} \slashsep \texttt{LLMs}}
	\entry
	{seit 9/2021}
	{IT Vorstand (ehrenamtlich)}
	{Feki.de e.V.}
	{\vspace{-10pt}
		\begin{itemize}[noitemsep,topsep=1pt,parsep=1pt,partopsep=0pt, leftmargin=-1pt]
			\item Betreuung und Weiterentwicklung der Vereinswebsite (\href{https://feki.de}{feki.de}, PHP-basiert, > 10.000 Besucher/Monat)
			\item Entwicklung der neuen Plattform (\href{https://new.feki.de}{new.feki.de}, Next.js + Strapi CMS)
			\item Leitung des vereinsinternen IT-Teams, zeitweise auch des offiziellen UniShops
			\item Serveradministration (u.a. Nextcloud, Gitlab, Mattermost, E-Mail-Server, auf zwei physischen Servern)
		\end{itemize}
		\texttt{Docker} \slashsep \texttt{CI/CD} \slashsep \texttt{Serveradministration} \slashsep \texttt{Next.js}}
	\entry
	{10/2019 -- 3/2023}
	{Studentische Hilfskraft (Übungsleiter)}
	{Lehrstuhl Praktische Informatik}
	{\vspace{-10pt}
		\begin{itemize}[noitemsep,topsep=1pt,parsep=1pt,partopsep=0pt, leftmargin=-1pt]
			\item Leitung von Übungsgruppen, teilweise auch Aufgabenkorrektur in den Kursen:
			\item Einführung in die Informatik: Algorithmen, Programmierung und Softwaretechnik (DSG-EiAPS-B)
			\item Java-Programmierung (DSG-JaP-B)
			\item Fortgeschrittene Java-Programmierung (DSG-AJP-B)
		\end{itemize}
		\texttt{C} \slashsep \texttt{Java} \slashsep \texttt{Formale Sprachen}}
\end{entrylist}

%----------------------------------------------------------------------------------------
%	EDUCATION
%----------------------------------------------------------------------------------------
\vspace{-10 pt}
\cvsect{Bildung}
\begin{entrylist}
	\entry
	{4/2023 -- 1/2026}
	{M. Sc. Angewandte Informatik}
	{Otto-Friedrich-Universität Bamberg}
	{Abschlussnote: Voraussichtlich 1,x (Bewertung der Masterarbeit ausstehend)}
	\entry
	{10/2019 -- 3/2023}
	{B. Sc. Angewandte Informatik}
	{Otto-Friedrich-Universität Bamberg}
	{Abschlussnote: 1,6 | Anwendungsgebiet: Statistik}
	\entry
	{9/2011 -- 6/2019}
	{Allgemeine Hochschulreife}
	{Gymnasium Münchberg}
	{Abschlussnote: 1,4}
\end{entrylist}

%----------------------------------------------------------------------------------------
%	EXPERIENCE
%----------------------------------------------------------------------------------------
\vspace{-10 pt}
\cvsect{Projekte}
\begin{entrylist}
	\entry
	{Masterarbeit}
	{Automatic Induction of Regular Constraints for a\\ Task-General Relation Extraction System}
	{\href{https://github.com/layaxx/masters-thesis}{github.com/layaxx/masters-thesis}}
	{Entwicklung eines Systems zur automatischen Induktion globaler Constraints, um die strukturelle Korrektheit in Sequenzmodellen zu erzwingen.\\
		\texttt{Python} \slashsep \texttt{Formale Sprachen} \slashsep \texttt{Transformer-Modelle}}
	\entry
	{Webanwendung}
	{Inventarverwaltungssystem für UniShop Bamberg}
	{\href{https://github.com/layaxx/unishop-inventory}{github.com/layaxx/unishop-inventory}}
	{Während meiner Leitungstätigkeit im offiziellen UniShop der Otto-Friedrich-Universität Bamberg entwickelt, unter anderem um PDF-Berichte über aktuelle Warenbestände zu generieren.\\
		\texttt{Next.js} \slashsep \texttt{TypeScript} \slashsep \texttt{LaTeX}}
	\entry
	{Webanwendung}
	{RSS-Feed Aggregator (Dashboard)}
	{\href{https://github.com/layaxx/dashboard}{github.com/layaxx/dashboard}}
	{Persönliches Dashboard zur Aggregation von RSS-Feeds und weiterer Hilfsprogramme.\\
		\texttt{Next.js} \slashsep \texttt{TypeScript} \slashsep \texttt{PostgreSQL}}
\end{entrylist}
%----------------------------------------------------------------------------------------
%	LANGUAGES
%----------------------------------------------------------------------------------------
\vspace{-10 pt}
\cvsect{Sprachen}
\vspace{-6pt}

\hspace{26mm} \textbf{Englisch} - fließend, \textbf{ Deutsch} - Muttersprache

%----------------------------------------------------------------------------------------

\end{document}
